\chapter[A short form of the third chapter's title]{A long long long long long long long third chapter title}
\label{cha:third-results}

\epigraph{%
    The ability to reduce everything to simple fundamental laws does not imply the ability to start from those laws and reconstruct the universe.}
    {P.W. Anderson~\cite{andersonMore1972}}
    %
    \begin{remarks}
        Meta-information about this chapter.
    \end{remarks}

    % The following space can be used as a "section abstract" where you summarize the work and its main findings 
    %
    \noindent \Ac{DFT} Lorem ipsum dolor sit amet, consetetur sadipscing elitr, sed diam nonumy eirmod tempor invidunt ut labore et dolore magna aliquyam erat, sed diam voluptua.
    %
    At vero eos et accusam et justo duo dolores et ea rebum.
    %
    Stet clita kasd gubergren, no sea takimata sanctus est Lorem ipsum dolor sit amet.
    %
    Lorem ipsum dolor sit amet, consetetur sadipscing elitr, sed diam nonumy eirmod tempor invidunt ut labore et dolore magna aliquyam erat, sed diam voluptua.
    %
    At vero eos et accusam et justo duo dolores et ea rebum.
    %
    Stet clita kasd gubergren, no sea takimata sanctus est Lorem ipsum dolor sit amet. \medskip
    %
%%%%%%%%%%%%%%%%%%%%%%%%%%%%%%%%%%%%%%%%%%%%%%%%%%%%%%%%%%%%%%%%%%%%

\section{Basics}
    \label{sec:third-basics}
    %
    \subsection{The ABC}
        \label{subsec:third-abc}
    \subsection{Counting to 10}
        \label{subsec:third-counting}

    \paragraph{How to add 1+1:} This works by adding 1+1, which gives 2.
    %
    \paragraph{How to multiply 2 by 2:} Do it, and find 4.

%
\section{Results}
    \label{sec:third-results}

%
\section{Discussion}
    \label{sec:third-discussion}